\documentclass[UTF8,11pt,a4paper]{ctexart}
\usepackage[margin=24mm]{geometry}
\usepackage{amsmath,amssymb,amsthm,mathtools,bm}
\usepackage{tikz-cd,enumitem,longtable,booktabs}
\usepackage{xurl}
\usepackage[unicode,hidelinks]{hyperref}
\allowdisplaybreaks
\setlength{\emergencystretch}{3em}
\setlength{\parindent}{0pt}
\setlength{\parskip}{0.45em}
\newtheorem*{theorem}{Theorem}
\newtheorem*{lemma}{Lemma}
\newtheorem*{proposition}{Proposition}
\newtheorem*{corollary}{Corollary}
\theoremstyle{definition}\newtheorem*{definition}{Definition}
\theoremstyle{remark}\newtheorem*{remark}{Remark}
\DeclareMathOperator{\Hom}{Hom}
\DeclareMathOperator{\End}{End}
\DeclareMathOperator{\Tr}{Tr}
\DeclareMathOperator{\rank}{rank}
\DeclareMathOperator{\Ind}{Ind}
\DeclareMathOperator{\Res}{Res}
\DeclareMathOperator{\im}{Im}
\DeclareMathOperator{\id}{id}
\newcommand{\R}{\mathbb R}
\newcommand{\C}{\mathbb C}
\newcommand{\Q}{\mathbb Q}
\newcommand{\Z}{\mathbb Z}
\newcommand{\N}{\mathbb N}
\newcommand{\E}{\mathbb E}
\newcommand{\Prob}{\mathbb P}
\newcommand{\eps}{\varepsilon}
\begin{document}
\section*{092｜有限群特征标的Artin诱导生成判据}
\textbf{作者／集体来源：}Pavel Etingof、Oleg Golberg、Sebastian Hensel、Tiankai Liu、Alex Schwendner、Elena Yudovina、Dmitry Vaintrob

\textbf{作品：}Introduction to Representation Theory，MIT 18.712，2010课程讲义，Chapter 4

\textbf{原文入口：}\url{https://ocw.mit.edu/courses/18-712-introduction-to-representation-theory-fall-2010/84358595a02a73bced2c4e363a5d66f0_MIT18_712F10_ch4.pdf}

\textbf{原文位置与计题范围：}PDF第8--10页的Definition 4.28、Theorems 4.32--4.33，以及第29--30页的Theorem 4.73和Corollary 4.74；只按Artin判据计一道。

\textbf{获取日期：}2026-09-28。\quad\textbf{复用依据：}MIT OCW，CC BY-NC-SA 4.0。使用2010课程PDF，不是2016出版书PDF。

\textbf{整理归属：}下列题面与证明取自所列人类来源，按注明范围转录、摘编；LaTeX环境、标题、换行及中文说明由助手整理。编辑调整在题内说明。

\section*{作者命题与人类证明}

\subsection*{Definition 4.28 and the construction of induction}
If $G$ is a group, $H\subset G$, and $V$ is a representation of $H$, then the induced representation $\Ind_H^G V$ is the representation of $G$ with
\[
\Ind_H^G V=\{f:G\to V\mid f(hx)=\rho_V(h)f(x)\text{ for all }x\in G,h\in H\}
\]
and the action $g(f)(x)=f(xg)$ for all $g\in G$.
In fact, $\Ind_H^G V$ is naturally isomorphic to $\Hom_H(k[G],V)$.
Let us check that $\Ind_H^G V$ is indeed a representation:
\[
g(f)(hx)=f(hxg)=\rho_V(h)f(xg)=\rho_V(h)g(f)(x),
\]
and
\[
g(g'(f))(x)=g'(f)(xg)=f(xgg')=(gg')(f)(x)
\]
for any $g,g',x\in G$ and $h\in H$.
Notice that if we choose a representative $x_\sigma$ from every right $H$-coset $\sigma$ of $G$, then any $f\in\Ind_H^G V$ is uniquely determined by $\{f(x_\sigma)\}$. Because of this,
\[
\dim(\Ind_H^G V)=\dim V\cdot\frac{|G|}{|H|}.
\]

\subsection*{Theorem 4.32: the Mackey formula}
Let us now compute the character $\chi$ of $\Ind_H^G V$. In each right coset $\sigma\in H\backslash G$, choose a representative $x_\sigma$.
\begin{theorem}[4.32]
One has
\[
\chi(g)=\sum_{\substack{\sigma\in H\backslash G\\x_\sigma gx_\sigma^{-1}\in H}}
\chi_V(x_\sigma gx_\sigma^{-1}).
\]
\end{theorem}
If the characteristic of the ground field $k$ is relatively prime to $|H|$, then this formula can be written as
\[
\chi(g)=\frac1{|H|}\sum_{\substack{x\in G\\xgx^{-1}\in H}}\chi_V(xgx^{-1}).
\]
\begin{proof}
For a right $H$-coset $\sigma$ of $G$, let us define
\[
V_\sigma=\{f\in\Ind_H^G V\mid f(g)=0\text{ for all }g\notin\sigma\}.
\]
Then one has $\Ind_H^G V=\bigoplus_\sigma V_\sigma$, and so
$\chi(g)=\sum_\sigma\chi_\sigma(g)$, where $\chi_\sigma(g)$ is the trace of the diagonal block of $\rho(g)$ corresponding to $V_\sigma$.
Since $g(\sigma)=\sigma g$ is a right $H$-coset for any right $H$-coset $\sigma$, $\chi_\sigma(g)=0$ if $\sigma\ne\sigma g$.
Now assume that $\sigma=\sigma g$. Then $x_\sigma g=hx_\sigma$ where $h=x_\sigma gx_\sigma^{-1}\in H$.
Consider the vector space homomorphism $\alpha:V_\sigma\to V$ with $\alpha(f)=f(x_\sigma)$. Since $f\in V_\sigma$ is uniquely determined by $f(x_\sigma)$, $\alpha$ is an isomorphism. We have
\begin{align*}
\alpha(gf)&=g(f)(x_\sigma)=f(x_\sigma g)=f(hx_\sigma)\\
&=\rho_V(h)f(x_\sigma)=h\alpha(f),
\end{align*}
and $gf=\alpha^{-1}h\alpha(f)$. This means that $\chi_\sigma(g)=\chi_V(h)$. Therefore
\[
\chi(g)=\sum_{\substack{\sigma\in H\backslash G\\\sigma g=\sigma}}\chi_V(x_\sigma gx_\sigma^{-1}).
\]
\end{proof}

\subsection*{Theorem 4.33: Frobenius reciprocity}
\begin{theorem}[4.33]
Let $H\subset G$ be groups, $V$ be a representation of $G$ and $W$ a representation of $H$. Then $\Hom_G(V,\Ind_H^G W)$ is naturally isomorphic to $\Hom_H(\Res_H^G V,W)$.
\end{theorem}
\begin{proof}
Let $E=\Hom_G(V,\Ind_H^G W)$ and $E'=\Hom_H(\Res_H^G V,W)$.
Define $F:E\to E'$ and $F':E'\to E$ as follows:
\[
F(\alpha)v=(\alpha v)(e),\qquad (F'(\beta)v)(x)=\beta(xv).
\]
In order to check that $F$ and $F'$ are well defined and inverse to each other, we need to check the following five statements. Let $\alpha\in E$, $\beta\in E'$, $v\in V$, and $x,g\in G$.

(a) $F(\alpha)$ is an $H$-homomorphism, i.e., $F(\alpha)hv=hF(\alpha)v$. Indeed,
\[
F(\alpha)hv=(\alpha hv)(e)=(h\alpha v)(e)=(\alpha v)(he)=(\alpha v)(eh)=h\cdot(\alpha v)(e)=hF(\alpha)v.
\]
(b) $F'(\beta)v\in\Ind_H^G W$. Indeed,
\[
(F'(\beta)v)(hx)=\beta(hxv)=h\beta(xv)=h(F'(\beta)v)(x).
\]
(c) $F'(\beta)$ is a $G$-homomorphism. Indeed,
\[
(F'(\beta)gv)(x)=\beta(xgv)=(F'(\beta)v)(xg)=(g(F'(\beta)v))(x).
\]
(d) $F\circ F'=\id_{E'}$. This holds since
\[
F(F'(\beta))v=(F'(\beta)v)(e)=\beta(v).
\]
(e) $F'\circ F=\id_E$. Indeed,
\[
(F'(F(\alpha))v)(x)=F(\alpha xv)=(\alpha xv)(e)=(x\alpha v)(e)=(\alpha v)(x),
\]
and we are done.
\end{proof}

\subsection*{Theorem 4.73: Artin's theorem (the counted problem)}
\begin{theorem}[4.73]
Let $X$ be a conjugation-invariant system of subgroups of a finite group $G$. Then two conditions are equivalent:

(i) Any element of $G$ belongs to a subgroup $H\in X$.

(ii) The character of any irreducible representation of $G$ belongs to the $\Q$-span of characters of induced representations $\Ind_H^G V$, where $H\in X$ and $V$ is an irreducible representation of $H$.
\end{theorem}
\begin{remark}
Statement (ii) of Theorem 4.73 is equivalent to the same statement with $\Q$-span replaced by $\C$-span. Indeed, consider the matrix whose columns consist of the coefficients of the decomposition of $\Ind_H^G V$ (for various $H,V$) with respect to the irreducible representations of $G$. Then both statements are equivalent to the condition that the rows of this matrix are linearly independent.
\end{remark}
\begin{proof}
\emph{Proof that (ii) implies (i).} Assume that $g\in G$ does not belong to any of the subgroups $H\in X$. Then, since $X$ is conjugation invariant, it cannot be conjugated into such a subgroup. Hence by the Mackey formula, $\chi_{\Ind_H^G(V)}(g)=0$ for all $H\in X$ and $V$. So by (ii), for any irreducible representation $W$ of $G$, $\chi_W(g)=0$. But irreducible characters span the space of class functions, so any class function vanishes on $g$, which is a contradiction.

\emph{Proof that (i) implies (ii).} Let $U$ be a virtual representation of $G$ over $\C$ (i.e., a linear combination of irreducible representations with nonzero integer coefficients) such that $(\chi_U,\chi_{\Ind_H^G V})=0$ for all $H,V$. So by Frobenius reciprocity, $(\chi_{U|H},\chi_V)=0$. This means that $\chi_U$ vanishes on $H$ for any $H\in X$. Hence by (i), $\chi_U$ is identically zero. This implies (ii) (because of the above remark).
\end{proof}
\begin{corollary}[4.74]
Any irreducible character of a finite group is a rational linear combination of induced characters from its cyclic subgroups.
\end{corollary}

\section*{整理说明与可修改标注（助手）}
\textbf{标准先修：}有限群复表示的完全可约性、不可约特征标的正交基性质、矩阵在有理数域与复数域上的秩一致。

\textbf{本题仍需完成的工作：}将子群对群元素的覆盖条件转化为诱导特征标的有理生成；须构造诱导表示、计算其迹并建立限制与诱导的Hom对应，而非仅引用Artin结论。

\textbf{取材说明：}保留作者定义、两项支撑定理及Artin主证明；题面指Theorem 4.73，不把支撑定理和Corollary 4.74另计题数。排版长等式只换行，原数学符号保留。

\textbf{$c$：}有限群特征标的有理生成

\textbf{$a$：}诱导表示的等变函数模型；陪集分块取迹；Frobenius互反；虚表示与整数重数矩阵。

\texttt{cross\_theory\_status = yes}

\textbf{具体联系及位置：}Theorem 4.73的两个方向分别把无法落入子群的元素变成特征标共同零点，并把正交补中的虚特征标限制到各子群；Theorems 4.32--4.33给出实现这两个转换的具体操作。

这些标注是非穷尽初稿；未列出不表示未使用。
\end{document}
